% ============================================================
% RESUMO / ABSTRACT
% ============================================================
\begin{abstract}
A interpretação automática do eletrocardiograma (ECG) é uma ferramenta
cada vez mais relevante para a triagem clínica, em especial para a
detecção tempestiva de condições cardíacas graves. Este trabalho propõe e
avalia arquiteturas de aprendizado profundo para a estimativa da
confiabilidade diagnóstica e a classificação multirrótulo de ECG na base
PTB-XL, sobre as cinco superclasses: ECG normal (NORM), infarto do
miocárdio (MI), alterações de ST/T (STTC), distúrbios de condução (CD) e
hipertrofia (HYP). As redes são supervisionadas diretamente pela
confiabilidade contínua anotada por cardiologistas, tratada como rótulo
suave (\textit{soft label}), e combinam um ramo temporal --- baseado em
convoluções dilatadas com blocos de \textit{squeeze-and-excitation} e
autoatenção --- a um ramo espectral derivado da transformada de Fourier.
Comparam-se duas variantes da mesma espinha dorsal: (A) um conjunto de
cinco redes especialistas independentes, uma por superclasse; e (B) um
único modelo de \textit{backbone} compartilhado com cinco saídas. Em ambas,
a presença de cada patologia é decidida por limiares otimizados por classe
no conjunto de validação e mantidos fixos no teste cego. A avaliação
abrange tanto a estimativa de confiança (MAE, RMSE) quanto a classificação
(AUC-ROC, AUPRC e F1 macro). No teste cego, ambas as variantes alcançaram
AUC-ROC macro de aproximadamente [XX,XX] e F1 macro de [XX,XX], com a
variante de \textit{backbone} compartilhado atingindo desempenho
equivalente ao das redes especialistas com cerca de um quinto dos
parâmetros. Os resultados indicam que a estimativa de confiabilidade
diagnóstica, aliada a uma decisão multirrótulo por limiares otimizados, é
um caminho promissor para sistemas auditáveis de apoio à decisão em
cardiologia.
\end{abstract}

\begin{IEEEkeywords}
eletrocardiograma, aprendizado profundo, classificação multirrótulo,
estimativa de confiança, PTB-XL, redes residuais
\end{IEEEkeywords}
